%% Companion technical report to the SANER 2027 submission (response plan section 13, 2026-09-13).
%%
%% WHAT THIS IS. The paper is self-contained on every claim; this report carries the procedures,
%% sensitivity analyses, and per-system tables the paper cites as "TR section N" (the \trX macros
%% in main.tex -- keep the SECTION ORDER here in sync with them). It is built from the SAME
%% generated/macros.tex and tables/*.tex as the paper, so no number is retyped, and it reproduces
%% every paper table so its captions' cross-references resolve.
%%
%% ANONYMITY. Shipped anonymized inside the reviewer mirror package (the CFP forbids a public
%% non-anonymized version during the review window); the named version with the real tool name
%% follows the 2026-12-01 notification under this distinct title. Same \ifanonymous switch as the
%% paper. BUILD: `make tr` (latexmk) in this directory.

\documentclass[10pt,conference]{IEEEtran}
\IEEEoverridecommandlockouts

\usepackage{cite}
\usepackage{amsmath,amssymb,amsfonts}
\usepackage{graphicx}
\usepackage{booktabs}
\usepackage{siunitx}
\usepackage{xcolor}
\usepackage{tikz}
\usepackage{url}
\usepackage[hidelinks]{hyperref}

\graphicspath{{figures/}}

\newif\ifanonymous
\anonymoustrue

\ifanonymous
  \newcommand{\tool}{\textsc{Anon}}
  \newcommand{\cli}{\texttt{arch}}
\else
  \newcommand{\tool}{\textsc{Anon}}
  \newcommand{\cli}{\texttt{arch}}
\fi
\newcommand{\TR}{this report}
\newcommand{\paper}{the paper}

%% GENERATED by analysis/make_macros.py -- DO NOT EDIT BY HAND.
%% Refreshing the data re-typesets the paper (plan §15).
\newcommand{\corpusCppCount}{5}
\newcommand{\corpusCsCount}{4}
\newcommand{\corpusReferenceCount}{9}
\newcommand{\corpusSystemCount}{9}
\newcommand{\rqAdequacyAdequateCount}{4}
\newcommand{\rqAdequacyCppLargestMax}{0.3028}
\newcommand{\rqAdequacyCppRatioMax}{0.1306}
\newcommand{\rqAdequacyFmtEntities}{66}
\newcommand{\rqAdequacyFmtLargest}{0.3182}
\newcommand{\rqAdequacyFmtRatio}{0.3333}
\newcommand{\rqAdequacyFmtUnits}{22}
\newcommand{\rqAdequacyInadequateCount}{5}
\newcommand{\rqAdequacyLargestIntervalHi}{0.7759}
\newcommand{\rqAdequacyLargestIntervalLo}{0.3028}
\newcommand{\rqAdequacyLargestMax}{0.5}
\newcommand{\rqAdequacyLibxmlLargestShare}{0.7759}
\newcommand{\rqAdequacyLibxmlRatio}{0.2414}
\newcommand{\rqAdequacyLooCorrect}{8/9}
\newcommand{\rqAdequacyLooLargestClauseFolds}{0}
\newcommand{\rqAdequacyLooWrong}{libxml2}
\newcommand{\rqAdequacyN}{9}
\newcommand{\rqAdequacyRatioIntervalHi}{1.0}
\newcommand{\rqAdequacyRatioIntervalLo}{0.1306}
\newcommand{\rqAdequacyRatioMax}{0.25}
\newcommand{\rqAdequacyRatioOnlyHiSystem}{libxml2}
\newcommand{\rqAdequacyRatioOnlyIntervalHi}{0.2414}
\newcommand{\rqAdequacyRhoBootReps}{2000}
\newcommand{\rqAdequacyRhoCiHi}{-0.093}
\newcommand{\rqAdequacyRhoCiLo}{-0.982}
\newcommand{\rqAdequacyRhoMqMoJo}{0.696}
\newcommand{\rqAdequacyRhoRatioAtoA}{-0.583}
\newcommand{\rqAdequacyRhoRatioMoJo}{-0.766}
\newcommand{\rqAdequacyRuleAgreement}{9/9}
\newcommand{\rqAdequacySerilogEntities}{6}
\newcommand{\rqAdequacySerilogLargest}{0.1667}
\newcommand{\rqAdequacySerilogRatio}{1.0}
\newcommand{\rqAdequacySerilogUnits}{6}
\newcommand{\rqAdequacyTerminalEntities}{63}
\newcommand{\rqAdequacyTerminalLargest}{0.0159}
\newcommand{\rqAdequacyTerminalRatio}{1.0}
\newcommand{\rqAdequacyTerminalUnits}{63}
\newcommand{\rqAdequacyTierBAdequateList}{none}
\newcommand{\rqAdequacyTierBInadequateList}{fmt, Serilog, Terminal}
\newcommand{\rqAdequacyTierBSystems}{3}
\newcommand{\rqFiveAcdcChurnMedian}{74.8}
\newcommand{\rqFiveAcdcChurnWorst}{47.7}
\newcommand{\rqFiveDirChurnMedian}{100.0}
\newcommand{\rqFivePipelineChurn}{100.0}
\newcommand{\rqNineEshopAri}{0.527}
\newcommand{\rqNineEshopBestSarAtoA}{85.87}
\newcommand{\rqNineEshopBestSarMoJo}{26.67}
\newcommand{\rqNineEshopBestSarName}{LIMBO}
\newcommand{\rqNineEshopConsensusClusters}{14.0}
\newcommand{\rqNineEshopCuratorAtoA}{95.79}
\newcommand{\rqNineEshopCuratorAtoASpread}{4.97}
\newcommand{\rqNineEshopCuratorClusters}{10.0}
\newcommand{\rqNineEshopCuratorCode}{R-A}
\newcommand{\rqNineEshopCuratorCount}{2}
\newcommand{\rqNineEshopCuratorEdgeFone}{0.629}
\newcommand{\rqNineEshopCuratorEdgeFoneSpread}{0.07}
\newcommand{\rqNineEshopCuratorFullRefAtoA}{TBD}
\newcommand{\rqNineEshopCuratorFullRefMoJo}{TBD}
\newcommand{\rqNineEshopCuratorMinusRatersMoJo}{20.15}
\newcommand{\rqNineEshopCuratorMoJo}{86.67}
\newcommand{\rqNineEshopCuratorMoJoSpread}{20.0}
\newcommand{\rqNineEshopCuratorPairAri}{0.747}
\newcommand{\rqNineEshopCuratorPairAtoA}{93.94}
\newcommand{\rqNineEshopCuratorPairKappaAligned}{0.824}
\newcommand{\rqNineEshopCuratorPairMoJo}{80.62}
\newcommand{\rqNineEshopCuratorPairNmi}{0.943}
\newcommand{\rqNineEshopCuratorVsConsensusAtoA}{92.0}
\newcommand{\rqNineEshopCuratorVsConsensusMoJo}{75.0}
\newcommand{\rqNineEshopFloorAtoA}{80.77}
\newcommand{\rqNineEshopFloorMoJo}{33.33}
\newcommand{\rqNineEshopGroups}{11.0}
\newcommand{\rqNineEshopKappa}{0.0}
\newcommand{\rqNineEshopKappaAligned}{0.702}
\newcommand{\rqNineEshopLabelsShared}{0}
\newcommand{\rqNineEshopLeakageDelta}{20.0}
\newcommand{\rqNineEshopMinutes}{43}
\newcommand{\rqNineEshopNmi}{0.868}
\newcommand{\rqNineEshopNonBlindMoJo}{66.67}
\newcommand{\rqNineEshopOps}{43.0}
\newcommand{\rqNineEshopOverBestSarMoJo}{60.0}
\newcommand{\rqNineEshopOverFloorMoJo}{53.34}
\newcommand{\rqNineEshopPairAgreement}{0.918}
\newcommand{\rqNineEshopProposeDerivedPct}{0}
\newcommand{\rqNineEshopRaterConflicts}{0}
\newcommand{\rqNineEshopRaterPairAtoA}{89.69}
\newcommand{\rqNineEshopRaterPairMoJo}{66.52}
\newcommand{\rqNineEshopRaterRefAriMin}{0.615}
\newcommand{\rqNineEshopRaterRefKappaAlignedMax}{0.937}
\newcommand{\rqNineEshopRaterRefKappaAlignedMin}{0.762}
\newcommand{\rqNineEshopRaterRefKappaMin}{0.0}
\newcommand{\rqNineEshopRaterRefMoJoMin}{74.16}
\newcommand{\rqNineEshopRaterRefNmiMin}{0.901}
\newcommand{\rqNineEshopRuleLines}{63.0}
\newcommand{\rqNineEshopSecondCuratorAtoA}{90.82}
\newcommand{\rqNineEshopSecondCuratorCode}{R-B}
\newcommand{\rqNineEshopSecondCuratorEdgeFone}{0.564}
\newcommand{\rqNineEshopSecondCuratorMinutes}{120}
\newcommand{\rqNineEshopSecondCuratorMoJo}{66.67}
\newcommand{\rqNineEshopSecondCuratorOps}{36.0}
\newcommand{\rqNineMedianCuratorMinusRatersMoJo}{17.8}
\newcommand{\rqNineSquidexAri}{0.228}
\newcommand{\rqNineSquidexBestSarAtoA}{92.96}
\newcommand{\rqNineSquidexBestSarMoJo}{60.0}
\newcommand{\rqNineSquidexBestSarName}{LIMBO}
\newcommand{\rqNineSquidexConsensusClusters}{12.0}
\newcommand{\rqNineSquidexCuratorAtoA}{95.83}
\newcommand{\rqNineSquidexCuratorClusters}{13.0}
\newcommand{\rqNineSquidexCuratorCode}{R-C}
\newcommand{\rqNineSquidexCuratorCount}{1}
\newcommand{\rqNineSquidexCuratorEdgeFone}{0.667}
\newcommand{\rqNineSquidexCuratorFullRefAtoA}{96.43}
\newcommand{\rqNineSquidexCuratorFullRefMoJo}{75.0}
\newcommand{\rqNineSquidexCuratorMinusRatersMoJo}{15.45}
\newcommand{\rqNineSquidexCuratorMoJo}{70.0}
\newcommand{\rqNineSquidexFloorAtoA}{92.0}
\newcommand{\rqNineSquidexFloorMoJo}{70.0}
\newcommand{\rqNineSquidexGroups}{12.0}
\newcommand{\rqNineSquidexKappa}{0.391}
\newcommand{\rqNineSquidexKappaAligned}{0.586}
\newcommand{\rqNineSquidexLabelsShared}{4}
\newcommand{\rqNineSquidexMinutes}{210}
\newcommand{\rqNineSquidexNmi}{0.766}
\newcommand{\rqNineSquidexOps}{46.0}
\newcommand{\rqNineSquidexOverBestSarMoJo}{10.0}
\newcommand{\rqNineSquidexOverFloorMoJo}{0.0}
\newcommand{\rqNineSquidexPairAgreement}{0.78}
\newcommand{\rqNineSquidexProposeDerivedPct}{0}
\newcommand{\rqNineSquidexRaterConflicts}{0}
\newcommand{\rqNineSquidexRaterPairAtoA}{86.3}
\newcommand{\rqNineSquidexRaterPairMoJo}{54.55}
\newcommand{\rqNineSquidexRaterRefAriMin}{n/a}
\newcommand{\rqNineSquidexRaterRefKappaAlignedMax}{n/a}
\newcommand{\rqNineSquidexRaterRefKappaAlignedMin}{n/a}
\newcommand{\rqNineSquidexRaterRefKappaMin}{n/a}
\newcommand{\rqNineSquidexRaterRefMoJoMin}{n/a}
\newcommand{\rqNineSquidexRaterRefNmiMin}{n/a}
\newcommand{\rqNineSquidexRuleLines}{59.0}
\newcommand{\rqNineSystems}{2}
\newcommand{\rqNineSystemsWithTwoCurators}{1}
\newcommand{\rqNineTwoCuratorSystemsList}{eShop}
\newcommand{\rqOneFloorAbseilAri}{44.06}
\newcommand{\rqOneFloorAbseilNested}{100.0}
\newcommand{\rqOneFloorAdequateAri}{64.88}
\newcommand{\rqOneFloorAdequateAtoA}{91.39}
\newcommand{\rqOneFloorAdequateMoJo}{94.4}
\newcommand{\rqOneFloorBashAri}{85.69}
\newcommand{\rqOneFloorCppAri}{44.06}
\newcommand{\rqOneFloorCppAtoA}{85.13}
\newcommand{\rqOneFloorCppMoJo}{91.94}
\newcommand{\rqOneFloorCsAri}{0.0}
\newcommand{\rqOneFloorCsAtoA}{80.48}
\newcommand{\rqOneFloorCsMoJo}{35.41}
\newcommand{\rqOneFloorItkAri}{22.61}
\newcommand{\rqOneFloorItkAtoA}{84.21}
\newcommand{\rqOneFloorItkNested}{96.15}
\newcommand{\rqOneFloorOpencvAri}{100.0}
\newcommand{\rqOneLlmEshopMoJo}{86.67}
\newcommand{\rqOneMedianAri}{84.07}
\newcommand{\rqOneMedianAtoA}{95.83}
\newcommand{\rqOneMedianMoJo}{72.11}
\newcommand{\rqOneProposeAbseilAri}{100.0}
\newcommand{\rqOneProposeAbseilMoJo}{100.0}
\newcommand{\rqOneProposeAdequateAri}{78.78}
\newcommand{\rqOneProposeCppMoJo}{82.66}
\newcommand{\rqOneProposeCsMoJo}{59.73}
\newcommand{\rqOneProposeEshopMoJo}{0.0}
\newcommand{\rqOneProposeItkAri}{84.6}
\newcommand{\rqOneProposeMedianMoJo}{69.47}
\newcommand{\rqOneProposeOpencvMoJo}{48.29}
\newcommand{\rqSevenColdBudget}{11/12}
\newcommand{\rqSevenItkClangSharePct}{1.5}
\newcommand{\rqSevenItkCmakeSharePct}{3.4}
\newcommand{\rqSevenItkColdMin}{410.0}
\newcommand{\rqSevenItkDoxygenSharePct}{93.4}
\newcommand{\rqSevenItkLZeroImpliedMin}{20.9}
\newcommand{\rqSevenItkWarmMin}{9.1}
\newcommand{\rqSevenMaxColdMin}{410.0}
\newcommand{\rqSevenMaxKloc}{825}
\newcommand{\rqSevenMaxPeakRssGb}{1.6}
\newcommand{\rqSevenMedianSpeedup}{4.28}
\newcommand{\rqSevenMinKloc}{0.1}
\newcommand{\rqSevenScalingSlope}{0.554}
\newcommand{\rqSevenSystems}{12}
\newcommand{\rqSevenWarmBudget}{10/12}
\newcommand{\rqSixAliasSurvival}{16/16}
\newcommand{\rqSixCombFindings}{458}
\newcommand{\rqSixCombFpTotal}{0}
\newcommand{\rqSixCombLoseLZeroAdvisory}{36}
\newcommand{\rqSixCombLoseLZeroDrift}{0}
\newcommand{\rqSixCombLoseLZeroGap}{422}
\newcommand{\rqSixCombLoseLZeroSilent}{0}
\newcommand{\rqSixCombModes}{4}
\newcommand{\rqSixCombMutations}{95}
\newcommand{\rqSixCombPrefixSilentOps}{add\_target, rename\_target}
\newcommand{\rqSixCombPrefixSilentTotal}{36}
\newcommand{\rqSixCombSilentModes}{none}
\newcommand{\rqSixCombSilentOps}{none}
\newcommand{\rqSixCombSilentTotal}{0}
\newcommand{\rqSixCombStarveMutatedAdvisory}{117}
\newcommand{\rqSixCombStarveMutatedDrift}{99}
\newcommand{\rqSixCombStarveMutatedGap}{242}
\newcommand{\rqSixCombStarveMutatedSilent}{0}
\newcommand{\rqSixCombStarveTenAdvisory}{31}
\newcommand{\rqSixCombStarveTenDrift}{404}
\newcommand{\rqSixCombStarveTenGap}{23}
\newcommand{\rqSixCombStarveTenSilent}{0}
\newcommand{\rqSixCombStarveThirtyAdvisory}{80}
\newcommand{\rqSixCombStarveThirtyDrift}{301}
\newcommand{\rqSixCombStarveThirtyGap}{77}
\newcommand{\rqSixCombStarveThirtySilent}{0}
\newcommand{\rqSixCombSystems}{5}
\newcommand{\rqSixDegradationSystems}{8}
\newcommand{\rqSixFOne}{1.0}
\newcommand{\rqSixFpPerPr}{0.0}
\newcommand{\rqSixHistArchitectural}{14}
\newcommand{\rqSixHistBuildCommits}{247}
\newcommand{\rqSixHistCommits}{1633}
\newcommand{\rqSixHistCoverageGaps}{0}
\newcommand{\rqSixHistFindings}{26}
\newcommand{\rqSixHistLabelsAgree}{25/26}
\newcommand{\rqSixHistLabelsKappa}{0.922}
\newcommand{\rqSixHistLabelsPending}{done}
\newcommand{\rqSixHistNewTargets}{9}
\newcommand{\rqSixHistNewUnmapped}{0}
\newcommand{\rqSixHistNoise}{0}
\newcommand{\rqSixHistNonArchitectural}{11}
\newcommand{\rqSixHistPairs}{8}
\newcommand{\rqSixHistPairsZero}{3}
\newcommand{\rqSixHistPairsZeroBuildChurn}{3}
\newcommand{\rqSixHistSystems}{3}
\newcommand{\rqSixMcNemarP}{\ensuremath{<}0.0001}
\newcommand{\rqSixMutationCount}{95}
\newcommand{\rqSixMutationSystems}{5}
\newcommand{\rqSixMutationsDetected}{95/95}
\newcommand{\rqSixOperatorCount}{6}
\newcommand{\rqSixPrecision}{1.0}
\newcommand{\rqSixRecall}{1.0}
\newcommand{\rqSixRenameAccuracy}{1.0}
\newcommand{\rqSixSuppressionRate}{1.0}
\newcommand{\rqSixWouldBeFalse}{5409}
\newcommand{\rqThreeAfiveEshopEdgesSaved}{30}
\newcommand{\rqThreeAfiveEshopRefRecallDrop}{1.0}
\newcommand{\rqThreeAfiveEshopRefRecallDropOperating}{0.0}
\newcommand{\rqThreeAfiveOperatingEdgesSaved}{0}
\newcommand{\rqThreeAfiveOperatingRankBiserial}{n/a}
\newcommand{\rqThreeAfiveOperatingSystemsChecked}{8}
\newcommand{\rqThreeAfiveOperatingSystemsMoved}{0}
\newcommand{\rqThreeAfiveRankBiserial}{1.0}
\newcommand{\rqThreeAfiveRefRecallDropList}{eShop 1.0, OrchardCore 0.41, nopCommerce 0.02}
\newcommand{\rqThreeAfiveRefRecallSystems}{3}
\newcommand{\rqThreeAfiveWilcoxonN}{7}
\newcommand{\rqThreeAfiveWilcoxonP}{0.0156}
\newcommand{\rqThreeCurationLiftCpp}{1.21}
\newcommand{\rqThreeCurationLiftCs}{33.86}
\newcommand{\rqThreeCurationRankBiserialAll}{1.0}
\newcommand{\rqThreeCurationRankBiserialCs}{1.0}
\newcommand{\rqThreeCurationWilcoxonP}{0.0312}
\newcommand{\rqThreeFloorCliffsCppVsCs}{0.6}
\newcommand{\rqThreeFloorCppMoJo}{92.31}
\newcommand{\rqThreeFloorCsMoJo}{35.41}
\newcommand{\rqThreeSevenMaxStructuralDelta}{0.0}
\newcommand{\rqThreeSevenSystemsChecked}{8}
\newcommand{\rqThreeSixLlmGtPipeline}{4}
\newcommand{\rqThreeSixLlmGtPipelineList}{abseil, eShop, libxml2, Squidex}
\newcommand{\rqThreeSixLlmMedianMoJo}{80.09}
\newcommand{\rqThreeSixLlmSystemsRan}{8}
\newcommand{\rqThreeSixPipelineGeLlm}{4/8}
\newcommand{\rqTwoAcdcMoJo}{28.12}
\newcommand{\rqTwoArcCsMoJo}{45.38}
\newcommand{\rqTwoArcMoJo}{36.92}
\newcommand{\rqTwoBaselineCount}{6}
\newcommand{\rqTwoClassicalAriMax}{78.76}
\newcommand{\rqTwoClassicalAriMaxWhere}{dir on bash}
\newcommand{\rqTwoCriticalDifference}{2.666}
\newcommand{\rqTwoCsAllCommonMedians}{LIMBO 40.88, WCA 39.76, dir 35.41, ACDC 15.44, cc 15.44}
\newcommand{\rqTwoCsBestAllCommonMoJo}{40.88}
\newcommand{\rqTwoCsBestAllCommonName}{LIMBO}
\newcommand{\rqTwoCsClassicalCells}{16}
\newcommand{\rqTwoCsClassicalFailed}{12}
\newcommand{\rqTwoCsClassicalPartial}{4}
\newcommand{\rqTwoCsClassicalPartialList}{nopCommerce: ARC 56.25; Squidex: WCA 50.0, LIMBO 60.0, ARC 50.0}
\newcommand{\rqTwoCsClassicalRecovered}{0}
\newcommand{\rqTwoFailedBelow}{50}
\newcommand{\rqTwoFloorAriWins}{3/9}
\newcommand{\rqTwoFloorAriWinsAll}{2/9}
\newcommand{\rqTwoFloorAriWinsOrTies}{3/9}
\newcommand{\rqTwoFloorAtoAWins}{2/9}
\newcommand{\rqTwoFloorAtoAWinsAll}{1/9}
\newcommand{\rqTwoFloorAtoAWinsOrTies}{2/9}
\newcommand{\rqTwoFloorCsRecovered}{0}
\newcommand{\rqTwoFloorMoJoWins}{2/9}
\newcommand{\rqTwoFloorMoJoWinsAll}{1/9}
\newcommand{\rqTwoFloorMoJoWinsOrTies}{5/9}
\newcommand{\rqTwoFloorNemenyiDiffers}{none}
\newcommand{\rqTwoFloorNemenyiNotDiffers}{ACDC, WCA, LIMBO, dir, cc}
\newcommand{\rqTwoFloorPosthocDiffers}{none}
\newcommand{\rqTwoFloorPosthocNotDiffers}{ACDC, WCA, LIMBO, dir, cc}
\newcommand{\rqTwoFriedmanAtoAP}{0.0095}
\newcommand{\rqTwoFriedmanFloorAtoAP}{0.1388}
\newcommand{\rqTwoFriedmanFloorN}{8}
\newcommand{\rqTwoFriedmanFloorP}{0.1205}
\newcommand{\rqTwoFriedmanFloorPPreExtension}{0.3484}
\newcommand{\rqTwoFriedmanN}{8}
\newcommand{\rqTwoFriedmanNPreExtension}{7}
\newcommand{\rqTwoFriedmanP}{0.0122}
\newcommand{\rqTwoFriedmanPPreExtension}{0.0491}
\newcommand{\rqTwoFriedmanProposeAtoAP}{0.15}
\newcommand{\rqTwoFriedmanProposeN}{8}
\newcommand{\rqTwoFriedmanProposeP}{0.1746}
\newcommand{\rqTwoLimboMoJo}{40.97}
\newcommand{\rqTwoLlmCostMaxUsd}{2.96}
\newcommand{\rqTwoLlmCostTotalUsd}{9.54}
\newcommand{\rqTwoLlmCsFailed}{OrchardCore}
\newcommand{\rqTwoLlmCsPartial}{nopCommerce, Squidex}
\newcommand{\rqTwoLlmCsRecovered}{eShop}
\newcommand{\rqTwoLlmMedianAri}{51.86}
\newcommand{\rqTwoLlmSystemsRan}{8}
\newcommand{\rqTwoNemenyiDiffers}{ACDC, cc}
\newcommand{\rqTwoNemenyiNotDiffers}{WCA, LIMBO, dir}
\newcommand{\rqTwoPipelineAriWins}{8/9}
\newcommand{\rqTwoPipelineAriWinsAll}{5/9}
\newcommand{\rqTwoPipelineAriWinsOrTies}{8/9}
\newcommand{\rqTwoPipelineAtoAWins}{7/9}
\newcommand{\rqTwoPipelineAtoAWinsAll}{4/9}
\newcommand{\rqTwoPipelineAtoAWinsOrTies}{7/9}
\newcommand{\rqTwoPipelineMoJo}{72.11}
\newcommand{\rqTwoPipelineMoJoWins}{6/9}
\newcommand{\rqTwoPipelineMoJoWinsAll}{4/9}
\newcommand{\rqTwoPipelineMoJoWinsOrTies}{7/9}
\newcommand{\rqTwoPipelineRank}{1.688}
\newcommand{\rqTwoPosthocDiffers}{ACDC, cc}
\newcommand{\rqTwoPosthocNotDiffers}{WCA, LIMBO, dir}
\newcommand{\rqTwoProposeAriWins}{5/9}
\newcommand{\rqTwoProposeAriWinsAll}{4/9}
\newcommand{\rqTwoProposeAriWinsOrTies}{5/9}
\newcommand{\rqTwoProposeAtoAWins}{4/9}
\newcommand{\rqTwoProposeAtoAWinsAll}{3/9}
\newcommand{\rqTwoProposeAtoAWinsOrTies}{4/9}
\newcommand{\rqTwoProposeMoJoWins}{4/9}
\newcommand{\rqTwoProposeMoJoWinsAll}{4/9}
\newcommand{\rqTwoProposeMoJoWinsOrTies}{4/9}
\newcommand{\rqTwoProposePosthocDiffers}{none}
\newcommand{\rqTwoProposePosthocNotDiffers}{ACDC, WCA, LIMBO, dir, cc}
\newcommand{\rqTwoRecoveredMin}{80}
\newcommand{\rqTwoSarAriMax}{32.59}
\newcommand{\rqTwoSarAriMaxWhere}{ACDC on OpenCV}
\newcommand{\rqTwoWcaMoJo}{32.03}
\newcommand{\tierPMedianAtoA}{100.0}
\newcommand{\tierPMedianFloorAtoA}{93.72}
\newcommand{\tierPMedianFloorMoJo}{96.87}
\newcommand{\tierPMedianMoJo}{100.0}
\newcommand{\tierPMedianProposeMoJo}{24.47}
\newcommand{\tierPSystemCount}{1}


\begin{document}

\ifanonymous
  \title{Build-Graph-First Architecture Recovery on Nine Open-Source Systems and One Industrial
  System: Procedures, Sensitivity Analyses, and Per-System Tables}
\else
  \title{Anon on Nine Open-Source Systems and One Industrial System: Procedures, Sensitivity
  Analyses, and Per-System Tables (Technical Report)}
\fi

\ifanonymous
  \author{\IEEEauthorblockN{Anonymous Author(s)}
  \IEEEauthorblockA{Affiliation withheld for double-anonymous review}}
\else
  \author{\IEEEauthorblockN{[anonymized] \textit{et al.}}
  \IEEEauthorblockA{(affiliations TBD)}}
\fi

\maketitle

\begin{abstract}
This technical report accompanies an empirical evaluation of a build-graph-first, deterministic
architecture-recovery pipeline for C/C++ and C\#\ systems. The paper is self-contained; this
report holds what its page budget could not: the per-build-idiom extraction rules, what every
baseline technique receives, the full inferential results, the sensitivity analysis of the
reference-free build-graph diagnostic, published scores on shared references, ablation and drift
detail per operator and evidence-loss mode, the per-system performance table, and the
rater-agreement and curation-effort detail of the blind-curation study. Every number here is a
generated macro or table from the same analysis pipeline as the paper; nothing is retyped.
\end{abstract}

\section{Scope and how to read this report}
\label{sec:scope}\label{sec:design}\label{sec:corpus}\label{sec:baselines}   % the shared terms paragraph and the table captions cite these; in the TR all point here

Section numbers below are the ones the paper cites (``TR \S$n$''). The terminology is the
paper's, restated below; the fidelity tables of the paper are reproduced in
\S\ref{sec:papertables} so that cross-references resolve. Nothing in this report is needed to
follow a claim in the paper; everything here is needed to reproduce or scrutinise one.

%% The shared terminology paragraph (was the boxed Table I until 2026-09-13; folded into prose
%% for the page budget, response plan section 13 status). The legend the table captions used to
%% carry lives here once; captions reference \S\ref{sec:design} (the enclosing section: IV in the
%% paper, 1 in the TR, which labels its scope section sec:design too) instead of restating it.
%% SHARED between main.tex and tr.tex: it must not reference a label that exists in only one of
%% the two documents. Hand-owned: definitions, not data.
\paragraph{Terminology.}
Three \emph{rows} recur in every table: the \emph{floor} (autonomous: every scored entity in its
owning build target, no curation), the \emph{propose} draft (autonomous: the tool's heuristic rule
draft accepted verbatim), and the \emph{curated} model (the committed hand rules applied). A
\emph{detection} number scores autonomous recovery, an \emph{expression} number what a human
expressed under evidence constraints; the two are never compared inferentially, and only detection
numbers support a recovery claim. Curation provenance (Cur.): \emph{sighted} = reference in view,
so the score is an oracle-assisted upper bound, never evidence against an autonomous technique;
\emph{assisted} = LLM-assisted from names and directories, the reference mapping not consulted;
\emph{blind} = reference-blind curator (RQ7); \emph{none} = curated model identical to the floor.
Reference grades (Ref) and the \emph{all-common} projection are defined in
\S\ref{sec:corpus} and \S\ref{sec:baselines}.


\section{Evidence extraction per build idiom}
\label{sec:tr-extract}

\begin{figure}[t]
  \centering
  \resizebox{\columnwidth}{!}{%% Hand-drawn (TikZ) example of ONE generated model, for the reviewer's "show a
%% representative generated architecture, traced to evidence" request. NOT generated by
%% analysis/ -- but every box and arrow below is transcribed 1:1 from a real pipeline
%% output, so re-run the pipeline before editing anything structural here:
%%
%%   subject   GNU bash 4.2 fixture (autotools L0 + Doxygen L2), run of 2026-07-16
%%   source    out/bash/architecture/generated/generated-components.dsl    (12 containers)
%%             out/bash/architecture/generated/generated-relationships.dsl (17 container->container lines)
%%             out/bash/architecture/generated/curated-facts.json          (evidence strings, verbatim)
%%   rules     overlays/bash/architecture/rules/mapping-rules.yaml (5 group: rules; 7 libs unmatched)
%%
%% Fidelity notes (keep these true):
%%  * Container names are the DSL names verbatim. All 17 container->container relationships
%%    of the DSL are drawn; nothing is added. The DSL's other 39 lines are component-altitude
%%    (a `bash`/`basename`/... component is one endpoint, e.g. `bash -> libbuiltins`); they
%%    belong to the component views, not to this container view, and are not drawn.
%%  * Line style = evidence strength from curated-facts.json (model.py compute_confidence):
%%      thick  = "link + include corroborate" (confidence medium)       3 edges
%%      solid  = "link only": declared in the Makefile, no include seen (confidence low, the
%%               plan's "declared but never used" case)                  6 edges
%%      dotted = "include only": not declared in any build file (low)   8 edges
%%  * Dashed boxes = the seven targets whose `tags` carry `needs-curation` in curated-facts
%%    (no mapping rule matched -> provisional singleton). The DSL additionally tags Shell
%%    Core / Example Loadables / Memory Allocator `needs-curation` through the rendering-time
%%    naming-convention lint (spaced names are the minority against seven lowercase lib
%%    names, generate_structurizr._mixed_convention_outliers); that lint is NOT drawn.
%%  * Evidence callouts 1-3 quote the `evidence[].detail` strings verbatim:
%%      1  rel:cpp:target:bash->cpp:target:libbuiltins
%%           link:    Makefile: bash links libbuiltins
%%           include: #include "builtins/common.h" (+16 more)     (count 17)
%%      2  rel:cpp:target:bash->cpp:target:libtermcap
%%           link:    Makefile: bash links libtermcap             (the only evidence)
%%      3  rel:cpp:target:libsh->cpp:target:bash
%%           include: #include "include/stdc.h" (+103 more)       (count 104)
%%  * No edge carries the `suspicious-declared` tag (architecture-drift.md: "_none_");
%%    the "link only" edges are the low-confidence declared-but-unused-in-source kind.
%%
%% Core tikz only (no \usetikzlibrary). Anonymity-safe: no tool name. Drawn at
%% ~9.6 x 6.9 cm for \resizebox{\columnwidth}{!}{...}; the caller wraps it in a figure.
\begin{tikzpicture}[x=1cm,y=1cm,
  % containers: rounded+shaded = curated C executable group; plain = C static library
  exe/.style={draw, rounded corners=2pt, fill=black!7, minimum height=0.5cm,
              inner xsep=3pt, inner ysep=3pt, align=center, font=\scriptsize\bfseries},
  lib/.style={draw, fill=white, minimum height=0.5cm, inner xsep=1.5pt, inner ysep=3pt,
              align=center, font=\scriptsize\ttfamily},
  nc/.style={dashed},                                   % needs-curation (no rule matched)
  % relationships, by evidence strength
  corr/.style={->, >=stealth, line width=1.4pt},        % link + include corroborate (medium)
  link/.style={->, >=stealth, thick},                   % declared link only (low)
  inc/.style={->, >=stealth, densely dotted, semithick},% include only, undeclared (low)
  % evidence callout markers + legend text
  ev/.style={circle, draw, fill=white, inner sep=0pt, minimum size=7.5pt,
             font=\tiny\bfseries},
  leg/.style={font=\scriptsize, inner sep=0pt, anchor=west},
  ]
  % ---- containers (names verbatim from generated-components.dsl) ------------------
  \node[exe] (intl) at (1.65,4.4) {Internationalization};
  \node[exe] (mem)  at (4.85,4.4) {Memory Allocator};
  \node[exe] (sup)  at (8.05,4.4) {Support \& Build Tools};
  \node[exe] (ex)   at (1.45,2.4) {Example Loadables};
  \node[exe] (core) at (4.85,2.4) [minimum width=2.6cm] {Shell Core};
  \node[lib, nc] (libsh) at (8.6,2.4) {libsh};
  % bottom row chained off each other's east anchors (cmtt widths differ per name)
  \node[lib, nc, anchor=west] (libbuiltins) at (0.15,0.4) {libbuiltins};
  \node[lib, nc, anchor=west] (libglob)     at ([xshift=1.8mm]libbuiltins.east) {libglob};
  \node[lib, nc, anchor=west] (libtilde)    at ([xshift=1.8mm]libglob.east)     {libtilde};
  \node[lib, nc, anchor=west] (libhistory)  at ([xshift=1.8mm]libtilde.east)    {libhistory};
  \node[lib, nc, anchor=west] (libreadline) at ([xshift=1.8mm]libhistory.east)  {libreadline};
  \node[lib, nc, anchor=west] (libtermcap)  at ([xshift=1.8mm]libreadline.east) {libtermcap};
  % ---- declared links from the shell core (Makefile LDADD, L0) ---------------------
  \draw[corr] ([xshift=-11mm]core.south) -- ([xshift=-1.5mm]libbuiltins.north)
    node[ev, pos=0.45, xshift=-4.5pt, yshift=4pt] {1};
  \draw[corr] ([xshift=-6mm]core.south)   -- ([xshift=-2mm]libglob.north);
  \draw[corr] ([xshift=-0.5mm]core.south) -- (libtilde.north);
  \draw[link] ([xshift=2.5mm]core.south)  -- (libhistory.north);
  \draw[link] ([xshift=5.5mm]core.south)  -- (libreadline.north);
  \draw[link] ([xshift=9mm]core.south)    -- (libtermcap.north)
    node[ev, pos=0.5, xshift=3.5pt, yshift=4.5pt] {2};
  \draw[link] ([yshift=1.4mm]core.east)   -- ([yshift=1.4mm]libsh.west);
  \draw[link] ([xshift=-8mm]core.north)   -- ([xshift=-4mm]intl.south east);
  % ---- include-only edges (Doxygen L2, no build-file declaration) ------------------
  \draw[inc] ([xshift=3mm]libbuiltins.north) -- ([xshift=-8.5mm]core.south);
  \draw[inc] ([xshift=2mm]libglob.north)       -- ([xshift=-3.5mm]core.south);
  \draw[inc] ([yshift=-1.4mm]libsh.west) -- ([yshift=-1.4mm]core.east)
    node[ev, pos=0.5, below=1pt] {3};
  \draw[inc] (ex.east) -- (core.west);
  \draw[inc] (ex.south) -- ([xshift=-5.5mm]libbuiltins.north);
  \draw[inc] (mem.south) -- (core.north);
  \draw[inc] ([xshift=5mm]core.north) -- ([xshift=-9mm]sup.south);
  \draw[inc] ([xshift=-5mm]sup.south) -- ([xshift=9mm]core.north);
  \draw[inc] (libsh.north) -- (libsh.north |- sup.south);
  % ---- legend: the three quoted evidence entries, then the line key ----------------
  \node[ev] (l1) at (0.30,-0.35) {1};
  \node[leg] at (l1.east) {\ \texttt{Makefile: bash links libbuiltins}\ +};
  \node[leg] at (0.55,-0.67) {\texttt{\#include "builtins/common.h" (+16 more)}};
  \node[ev] (l2) at (0.30,-0.99) {2};
  \node[leg] at (l2.east) {\ \texttt{Makefile: bash links libtermcap}};
  \node[ev] (l3) at (0.30,-1.31) {3};
  \node[leg] at (l3.east) {\ \texttt{\#include "include/stdc.h" (+103 more)}};
  \draw[corr] (0.15,-1.70) -- (0.65,-1.70) node[leg, xshift=3pt] (k1) {link + include (medium)};
  \draw[link] ([xshift=10pt]k1.east) -- ++(0.5,0) node[leg, xshift=3pt] (k2) {declared link only (low)};
  \draw[inc]  ([xshift=10pt]k2.east) -- ++(0.5,0) node[leg, xshift=3pt] (k3) {include only (low)};
  \draw[nc]   (0.15,-2.15) rectangle (0.65,-1.89);
  \node[leg] at (0.76,-2.02) {needs-curation: no mapping rule matched this target (provisional singleton)};
\end{tikzpicture}
}
  \caption{One generated model with its evidence: the container view generated for GNU bash 4.2,
  transcribed from the emitted DSL. Twelve containers and all seventeen container-level
  relationships; line weight shows evidence strength (thick: declared link corroborated by
  includes; solid: declared only; dotted: includes without a declaration), dashed boxes are
  targets no rule claimed. Every arrow cites evidence; callouts 1--3 quote three verbatim,
  including a Makefile link no source file exercises (2).}
  \label{fig:model}
\end{figure}

The pipeline distinguishes evaluated from static extraction, and the difference bounds what a
consistently extracted graph can miss. Figure~\ref{fig:model} shows what the extracted evidence
becomes after curation and generation: every arrow of the generated bash model is traceable to
the build-file line or include that produced it.

\paragraph{CMake (evaluated).} The project is configured once with the Ninja generator and
default options unless a preset or flags are given; the per-system flags are in the artifact. The
pipeline then reads CMake's File API codemodel of that configured tree: targets of every type
(executables, static, shared, object, and interface libraries; utility targets typed unknown), the
link dependencies as CMake evaluated them for that configuration, generator expressions and
conditions included, and the \texttt{PUBLIC}/\texttt{PRIVATE} visibility of each edge from CMake's
own graphviz output, against which the codemodel edges are cross-checked. Alias targets collapse
onto their real target, imported and external targets are excluded, and a target whose generated
sources are absent is tagged partial.

\paragraph{.NET (static L0, evaluated L2).} The L0 layer is a static parse of \texttt{.sln} and
\texttt{.csproj} files that takes every \texttt{<ProjectReference>} regardless of MSBuild
conditions and follows no import; the Roslyn layer then loads the solution through
\texttt{MSBuildWorkspace}, an evaluated view, and contributes symbol references and the per-document
owning project. A project that fails to restore contributes no symbol edges: the graph is then
honestly partial and the coverage record says so.

\paragraph{MSBuild C++ and automake (static).} \texttt{.vcxproj} files and automake or
hand-written \texttt{Makefile.in} are parsed statically: imports, property expansion, and
conditionals are not evaluated, so a conditional reference is over-approximated and one reachable
only through an import is missed. The automake model also supplies target ownership for Doxygen
attribution when CMake is absent.

\paragraph{Layers and coverage.} Declared dependencies form evidence layer L0; package references
form L1; L2 adds the \texttt{\#include} graph (Doxygen XML or \texttt{clang-scan-deps}) and Roslyn
symbol references; L3 is reserved for fine-grained evidence (a P/Invoke interop extractor and a
deployment-descriptor extractor exist but are not evaluated). Every extractor writes an independent
JSON fragment and degrades gracefully: without its toolchain it records a coverage warning and an
empty fragment. Per-layer coverage travels with the facts; its L0 denominator is the first-party
target set, full or near-full on every corpus system (per run in the artifact). The extracted
graph is not validated against an independent evaluated source (paper, \S{}VII); a spot check
against \texttt{cmake --graphviz} is future work.

\section{Baseline inputs, entity ownership, and the LLM runner}
\label{sec:tr-inputs}

One dependency graph per system, exported once as RSF at the scored granularity, feeds every
technique (paper, \S{}IV-B). This section records the assignment rules the paper summarises.

\paragraph{Owning build target.} The pipeline's autonomous rows know each entity's owning build
target, which no baseline receives in any form. A file belongs to the CMake target whose source
list names it, else to the target owning the longest directory prefix (Makefile ownership
analogously); a file two targets list goes to the later one in the codemodel; files no target
claims stay unclustered and count against the floor. For ITK the unclustered files are mostly
headers outside any source list. Per-system counts of clustered and unclustered entities are in
the artifact's graph sidecars (\texttt{generated/graph/}).

\paragraph{What each technique receives.} At project granularity (C\#) the graph holds every
relationship kind between first-party projects (project references, symbol uses, interop) as one
unweighted edge set; at file granularity (C++) it holds the \texttt{\#include} graph, and declared
link edges between targets are never projected onto files. ACDC and \texttt{cc} see the edges
only; \texttt{dir} sees entity paths; WCA and LIMBO also receive $k$, the reference cluster count,
as their stopping point; ARC receives $k$ and each entity's source text; the LLM receives the
graph, one line of repository metadata, and $k$.

\paragraph{LLM-recovery runner.} A reasoning-tier model (\texttt{gpt-5.6-sol} on Azure AI
Foundry, pinned by identifier) receives the exported graph, the metadata line, and $k$ under one
prompt asking for a JSON entity-to-cluster map, at default sampling (the model rejects a
temperature override), under a 180\,k input-token budget above which the runner refuses (ITK
exceeds it). Five reruns per system on \rqTwoLlmSystemsRan{} systems, at most
\$\rqTwoLlmCostMaxUsd{} per system at list price, \$\rqTwoLlmCostTotalUsd{} in total; the mean is
reported in the paper's Table~\ref{tab:rq2} with the rerun standard deviation, and the prompt,
every reply, and the per-rerun cluster counts are in the artifact.

\section{Metrics and inferential procedures}
\label{sec:tr-stats}

\paragraph{Partition metrics.} MoJoFM~\cite{mojofm} is the normalised move-and-join distance;
a2a~\cite{le2015a2a} additionally charges for adding, removing, and moving clusters; ARI is the
chance-corrected adjusted Rand index ($\times$100), added post hoc at a reviewer's request because
MoJoFM's cheap joins reward over-split partitions and a2a's range is
compressed~\cite{zhang2023sarif}. Beside ARI we print the \emph{nested share}, the fraction of a
partition's clusters lying entirely inside one reference cluster, which tells over-splitting from
misgrouping. c2c and the adjusted a2a of Zhang et al.\ are not reported (the latter has no released
implementation). \emph{Container-edge F1} scores the lifted container edges against the reference's
edge set, the dependency graph projected through the reference's own clusters after a greedy
maximum-overlap label alignment.

\paragraph{Inter-rater agreement.} Reported label-free, since independent raters name clusters
differently and a raw label-match $\kappa$ then measures wording: ARI, normalized mutual
information (NMI), and Cohen's $\kappa$ after optimal (Hungarian) label matching, with the raw
$\kappa$ beside them because the plan's flag rule (raw $\kappa<0.6$ marks a low-confidence
reference) is applied as written.

\paragraph{RQ2 inferential tests (autonomous rows only).} A Friedman omnibus on MoJoFM and a2a
over the all-common techniques (ACDC, WCA, LIMBO, \texttt{dir}, \texttt{cc}) with the floor in
place of the curated row does not reject equality: MoJoFM $p=\rqTwoFriedmanFloorP$, a2a
$p=\rqTwoFriedmanFloorAtoAP$, $n=\rqTwoFriedmanFloorN$ systems on which every technique completed,
ARC excluded. Before Squidex was admitted to the corpus the same omnibus gave
$p=\rqTwoFriedmanFloorPPreExtension$. With the propose draft in place of the curated row:
$p=\rqTwoFriedmanProposeP$ (MoJoFM) and \rqTwoFriedmanProposeAtoAP{} (a2a). Post hoc, neither row
differs from any technique under Holm-corrected pairwise Wilcoxon signed-rank tests or under
Nemenyi with critical difference \rqTwoCriticalDifference{}. The curated row enters no inferential
test; its omnibus, $p=\rqTwoFriedmanP$, differing from \rqTwoPosthocDiffers{} after correction,
is descriptive of assisted against unassisted recovery, since three of its C\#\ members saw the
reference.

\paragraph{Other tests.} Paired on/off comparisons (RQ3) use Wilcoxon signed-rank with the
matched-pairs rank-biserial $r$; Cliff's $\delta$ is reserved for the one between-group contrast
(C++ vs.\ C\#\ floor); detection rates get Wilson intervals, the stopping rules exact McNemar, and
the adequacy diagnostic Spearman's $\rho$ with a bootstrap interval over
\rqAdequacyRhoBootReps{} resamples. With $N\le 6$ paired systems the signed-rank $p$ cannot fall
below $0.0625$ however large the effect. Every statistic is a generated macro from the raw run
outputs (\texttt{analysis/stats.py}).

\section{The reference-free diagnostic: sensitivity}
\label{sec:tr-diagnostic}

\begin{figure}[t]
  \centering
  \resizebox{\columnwidth}{!}{% GENERATED by analysis/make_figures.py from stats.json -- DO NOT EDIT BY HAND.
\begin{tikzpicture}
  \draw[thick] (0,0) rectangle (6.4,4);
  \draw (0.963,0) -- (0.963,-0.09); \node[below,font=\footnotesize] at (0.963,-0.09) {0.02};
  \draw (2.237,0) -- (2.237,-0.09); \node[below,font=\footnotesize] at (2.237,-0.09) {0.05};
  \draw (3.2,0) -- (3.2,-0.09); \node[below,font=\footnotesize] at (3.2,-0.09) {0.1};
  \draw (4.163,0) -- (4.163,-0.09); \node[below,font=\footnotesize] at (4.163,-0.09) {0.2};
  \draw (5.437,0) -- (5.437,-0.09); \node[below,font=\footnotesize] at (5.437,-0.09) {0.5};
  \draw (6.4,0) -- (6.4,-0.09); \node[below,font=\footnotesize] at (6.4,-0.09) {1};
  \draw (0,0) -- (-0.09,0); \node[left,font=\footnotesize] at (-0.09,0) {0};
  \draw (0,0.8) -- (-0.09,0.8); \node[left,font=\footnotesize] at (-0.09,0.8) {20};
  \draw (0,1.6) -- (-0.09,1.6); \node[left,font=\footnotesize] at (-0.09,1.6) {40};
  \draw (0,2.4) -- (-0.09,2.4); \node[left,font=\footnotesize] at (-0.09,2.4) {60};
  \draw (0,3.2) -- (-0.09,3.2); \node[left,font=\footnotesize] at (-0.09,3.2) {80};
  \draw (0,4) -- (-0.09,4); \node[left,font=\footnotesize] at (-0.09,4) {100};
  \node[below,font=\small] at (3.2,-0.5) {unit ratio (build units / scored entities, log)};
  \node[rotate=90,font=\small] at (-0.95,2) {floor MoJoFM (all-common)};
  \draw[densely dashed,black!60] (4.473,0) -- (4.473,4);
  \node[anchor=south east,rotate=90,font=\tiny,black!60] at (4.473,1.2) {ratio $\le$ 0.25};
  \draw[densely dashed,black!60] (0,3.2) -- (6.4,3.2);
  \node[anchor=south west,font=\tiny,black!60] at (0.132,3.2) {recovered $\ge$ 80};
  \node[font=\footnotesize,black!55] at (2.237,2.48) {build-modular};
  \node[font=\footnotesize,black!55] at (5.29,2) {project-coarse};
  \fill (3.435,3.597) circle (2pt);
  \node[anchor=east,font=\tiny] at (3.365,3.597) {abseil};
  \fill (3.571,3.678) circle (2pt);
  \node[anchor=west,font=\tiny] at (3.641,3.618) {bash};
  \draw (6.4,1.333) circle (2pt);
  \node[anchor=east,font=\tiny] at (6.32,1.263) {eShop};
  \fill (0.526,3.874) circle (2pt);
  \node[anchor=west,font=\tiny] at (0.596,3.744) {ITK};
  \draw (4.425,0.457) circle (2pt);
  \node[anchor=east,font=\tiny] at (4.355,0.457) {libxml2 (largest unit 78\,\%)};
  \draw (6.4,1.5) circle (2pt);
  \node[anchor=east,font=\tiny] at (6.32,1.6) {nopCommerce};
  \fill (0.468,4) circle (2pt);
  \node[anchor=west,font=\tiny] at (0.538,3.96) {OpenCV};
  \draw (6.4,0.463) circle (2pt);
  \node[anchor=east,font=\tiny] at (6.32,0.463) {OrchardCore};
  \draw (6.4,2.8) circle (2pt);
  \node[anchor=east,font=\tiny] at (6.32,2.8) {Squidex};
\end{tikzpicture}
}
  \caption{The paper's two-regime figure: the autonomous floor's MoJoFM against the build graph's
  unit ratio. Dashed: the post-hoc rule's ratio threshold (\rqAdequacyRatioMax{}) and the recovery
  threshold (\rqTwoRecoveredMin{}). Filled markers satisfy the rule; libxml2 fails it through the
  largest-unit-share clause.}
  \label{fig:regimes}
\end{figure}

The floor's \emph{unit ratio} counts the build units that own at least one scored entity over the
scored entities (a target owning no scored file does not count, which is why the ratio lies below
the target counts of Table~\ref{tab:rq1}); its \emph{largest-unit share} is the fraction of scored
entities in the largest unit; TurboMQ is computed over the exported graph. Across the
\rqAdequacyN{} scored systems the unit ratio correlates with the floor's MoJoFM at Spearman
$\rho=\rqAdequacyRhoRatioMoJo{}$ (bootstrap 95\,\% interval
$[\rqAdequacyRhoCiLo{}, \rqAdequacyRhoCiHi{}]$ over \rqAdequacyRhoBootReps{} resamples;
$\rqAdequacyRhoRatioAtoA{}$ with a2a) and TurboMQ at $\rho=\rqAdequacyRhoMqMoJo{}$.

The post-hoc rule ``unit ratio $\le\rqAdequacyRatioMax{}$ and largest-unit share
$\le\rqAdequacyLargestMax{}$'' declares \rqAdequacyAdequateCount{} systems build-modular (ITK,
abseil, OpenCV, bash: ratios up to \rqAdequacyCppRatioMax{}, largest unit up to
\rqAdequacyCppLargestMax{}) and \rqAdequacyInadequateCount{} not (every C\#\ system at ratio 1.0;
libxml2, whose one library holds \rqAdequacyLibxmlLargestShare{} of the files), and agrees with a
floor MoJoFM of at least \rqTwoRecoveredMin{} on \rqAdequacyRuleAgreement{} systems
(Fig.~\ref{fig:regimes}).

\paragraph{Threshold intervals.} The ratio threshold can sit anywhere in
$[\rqAdequacyRatioIntervalLo{}, \rqAdequacyRatioIntervalHi{})$ without changing any call. Without
the largest-unit clause, the ratio alone separates the corpus only inside
$[\rqAdequacyRatioIntervalLo{}, \rqAdequacyRatioOnlyIntervalHi{})$: libxml2 sits at a low ratio
with one dominant library, which is what the second clause exists for.

\paragraph{Leave-one-system-out.} Re-selecting both thresholds with each system held out predicts
\rqAdequacyLooCorrect{} systems correctly and misses \rqAdequacyLooWrong{}, the one system the
largest-unit clause exists for: no training fold selects that clause, so the clause is a
one-system clause and the rule is a candidate diagnostic, not a validated predictor. The rule was
chosen after seeing these systems, which co-vary in language, build idiom, granularity, and
reference provenance.

\paragraph{Tier-B calls.} Of the \rqAdequacyTierBSystems{} unreferenced Tier-B systems the rule
calls none build-modular: Serilog and Terminal sit at ratio 1.0, and fmt (\rqAdequacyFmtUnits{}
CMake targets over \rqAdequacyFmtEntities{} files, ratio \rqAdequacyFmtRatio{}) is the
project-coarse C++ case the scored corpus lacks. These calls are predictions the extension corpus
can test once references exist.

\section{Published scores on shared references}
\label{sec:tr-published}

SARIF~\cite{zhang2023sarif} scored bash 4.2, ITK 4.5.2, and libxml2 2.4.22 against the references
and versions used in the paper: the libxml2 reference is SARIF's own industry-labelled ground
truth, and bash and ITK are the Lutellier et al.\ ground truths SARIF also scores against. Its
published MoJoFM/a2a and those of FCA~\cite{teymourian2022fca}, which SARIF ran as a baseline, can
therefore be read beside Table~\ref{tab:rq2}, descriptively only, since they score the authors'
own entity set rather than our all-common projection:

\begin{center}
\footnotesize\setlength{\tabcolsep}{3.5pt}
\begin{tabular}{@{}l rr rr rr@{}}
\toprule
 & \multicolumn{2}{c}{bash} & \multicolumn{2}{c}{ITK} & \multicolumn{2}{c}{libxml2} \\
\cmidrule(lr){2-3}\cmidrule(lr){4-5}\cmidrule(lr){6-7}
Published (own entity set) & MoJo & a2a & MoJo & a2a & MoJo & a2a \\
\midrule
SARIF~\cite{zhang2023sarif}, Table 4 & 70 & 87 & 68 & 86 & 49 & 89 \\
FCA, as run in~\cite{zhang2023sarif} & 41 & 80 & 44 & 73 & 37 & 84 \\
ACDC, as run in~\cite{zhang2023sarif} & 50 & 81 & 40 & 73 & -- & -- \\
\bottomrule
\end{tabular}
\end{center}

On bash and ITK our autonomous floor sits above both SARIF and FCA, and SARIF's own ACDC run lies
below our ACDC row, so our inputs do not handicap the classical baselines. On libxml2 SARIF's
fusion of code text and folder structure beats every structural partition we score, and only the
LLM baseline scores higher: one library target gives the build graph nothing to report there.
SARIF's only release is a closed Linux binary for C/C++/Java via Depends, without C\#\ support and
without a licence; FCA exists as one unlicensed MATLAB script. Neither was run.

\section{Ablation detail}
\label{sec:tr-ablation}

\paragraph{Evidence layers (A1).} Capping the ladder at L0 already yields most container edges;
the symbol layer (L2) adds the remainder and L1 (package references) none at container altitude.
The post-hoc layer filter agrees with the engine's own evidence ceiling on all
\rqThreeSevenSystemsChecked{} systems, i.e.\ no container edge exists whose only evidence is a
layer that was capped away. The per-layer container-edge counts are the A1 columns of
Table~\ref{tab:rq1}.

\paragraph{Weight threshold (A3) and visibility down-weighting (A4).} Sweeping the significance
threshold barely moves the edge set at the operating point (3), because declared edges clear it on
their base weight and the include layer's cap of 10 keeps chatty headers from dominating; A4
(down-weighting \texttt{PRIVATE} link edges) is inert on every system, since visibility changes
weight but no edge crosses the threshold as a result.

\paragraph{Never-drop-declared (A5).} At the default threshold the rule saves
\rqThreeAfiveOperatingEdgesSaved{} edges on \rqThreeAfiveOperatingSystemsChecked{} systems. At
threshold 8 a scalar filter removes \rqThreeAfiveEshopEdgesSaved{} of eShop's container edges;
edge counts move on \rqThreeAfiveWilcoxonN{} systems (paired $r=\rqThreeAfiveRankBiserial{}$,
$p=\rqThreeAfiveWilcoxonP{}$), and container-edge recall against the reference falls by
\rqThreeAfiveRefRecallDropList{} on the \rqThreeAfiveRefRecallSystems{} systems where it is scored.

\paragraph{LLM suggestions and naming (A6, A7).} Handed the same graph, the LLM baseline beats the
curated model on \rqThreeSixLlmGtPipelineList{}; the deterministic pipeline matches or beats it on
\rqThreeSixPipelineGeLlm{} of the \rqThreeSixLlmSystemsRan{} systems it completed (median LLM
MoJoFM \rqThreeSixLlmMedianMoJo{}); its recovery varies across reruns on Squidex (Table~\ref{tab:rq2}
gives the dispersion), and ITK's file graph exceeded the single-prompt input budget, so no call was made. An adversarial naming provider that renames every element
perturbs \rqThreeSevenMaxStructuralDelta{} relationships on \rqThreeSevenSystemsChecked{} systems.

\section{Drift arms}
\label{sec:tr-drift}
%% GENERATED by analysis/make_tables.py from results.csv -- DO NOT EDIT BY HAND.
\begin{table}[t]
  \centering
  \caption{RQ5 seeded-mutation detection per operator, pooled over 5 systems (first-party L0 scope): Mut.\ = injected mutations, Det.\ = mutations every expected finding of which was reported, TP/FN/FP = drift-report findings (one per vanished or new target, edge, rename, or violation, so a removed target yields several).}
  \label{tab:rq6}
  \footnotesize\setlength{\tabcolsep}{3.5pt}
  \begin{tabular}{@{}l rrrrr rrr@{}}
    \toprule
    Operator & Mut. & Det. & TP & FN & FP & Prec. & Rec. & F1 \\
    \midrule
    \texttt{add\_forbidden\_dep} & 12 & 12 & 24 & 0 & 0 & 100.0 & 100.0 & 100.0 \\
    \texttt{remove\_dep} & 16 & 16 & 16 & 0 & 0 & 100.0 & 100.0 & 100.0 \\
    \texttt{add\_target} & 20 & 20 & 20 & 0 & 0 & 100.0 & 100.0 & 100.0 \\
    \texttt{remove\_target} & 20 & 20 & 98 & 0 & 0 & 100.0 & 100.0 & 100.0 \\
    \texttt{rename\_target} & 16 & 16 & 196 & 0 & 0 & 100.0 & 100.0 & 100.0 \\
    \texttt{merge\_targets} & 11 & 11 & 104 & 0 & 0 & 100.0 & 100.0 & 100.0 \\
    \midrule
    \textbf{overall} & 95 & 95 & 458 & 0 & 0 & \textbf{100.0} & \textbf{100.0} & \textbf{100.0} \\
    \bottomrule
  \end{tabular}
\end{table}

%% GENERATED by analysis/make_tables.py from results.csv -- DO NOT EDIT BY HAND.
\begin{table}[t]
  \centering
  \caption{RQ5: seeded mutations replayed with evidence loss induced on the mutated run, pooled over 5 systems -- where each injected finding surfaced. Drift = reported as named drift; Gap = a labelled coverage gap; Adv. = visible only through the run's advisory downgrade; Silent = hidden.}
  \label{tab:rq6b}
  \footnotesize\setlength{\tabcolsep}{3.5pt}
  \begin{tabular}{@{}l rrrrrrr@{}}
    \toprule
    Evidence loss & Mut. & Find. & Drift & Gap & Adv. & Silent & FP \\
    \midrule
    10\% targets starved & 95 & 458 & 404 & 23 & 31 & 0 & 0 \\
    30\% starved (floor) & 95 & 458 & 301 & 77 & 80 & 0 & 0 \\
    mutated targets starved & 95 & 458 & 99 & 242 & 117 & 0 & 0 \\
    L0 extractor lost & 95 & 458 & 0 & 422 & 36 & 0 & 0 \\
    \bottomrule
  \end{tabular}
\end{table}

%% GENERATED by analysis/make_tables.py from results.csv -- DO NOT EDIT BY HAND.
\begin{table}[t]
  \centering
  \caption{RQ5: real release history, fixed consecutive releases (not selected by changelog), rules held constant across each pair. Build = commits touching build files; New = first-party targets the later release adds that the rules leave unmapped; Labels = findings both raters labelled identically (architectural/non-architectural/noise; -- = pending; disagreements and $\kappa$ in the text).}
  \label{tab:rq6hist}
  \footnotesize\setlength{\tabcolsep}{3pt}
  \begin{tabular}{@{}l rrrrrl@{}}
    \toprule
    Release pair & Commits & Build & Find. & Gaps & New & Labels \\
    \midrule
    libxml2 v2.13.0$\to$v2.14.0 & 673 & 108 & 2 & 0 & 0 & 0/2/0 \\
    libxml2 v2.14.0$\to$v2.15.0 & 368 & 51 & 0 & 0 & 0 & 0/0/0 \\
    libxml2 v2.15.0$\to$v2.15.3 & 79 & 10 & 2 & 0 & 0 & 0/2/0 \\
    nopCommerce 4.70.0$\to$4.80.0 & 398 & 49 & 12 & 0 & 0 & 8/4/0 \\
    Squidex 7.17.0$\to$7.18.0 & 16 & 6 & 8 & 0 & 0 & 5/3/0 \\
    Squidex 7.18.0$\to$7.19.0 & 29 & 7 & 2 & 0 & 0 & 1/0/0 \\
    Squidex 7.20.0$\to$7.21.0 & 19 & 5 & 0 & 0 & 0 & 0/0/0 \\
    Squidex 7.21.0$\to$7.22.0 & 51 & 11 & 0 & 0 & 0 & 0/0/0 \\
    \bottomrule
  \end{tabular}
\end{table}


\paragraph{Seeded mutations (Table~\ref{tab:rq6}).} \rqSixMutationCount{} labelled changes over
\rqSixOperatorCount{} operators were injected into the build files of checkout copies of
\rqSixMutationSystems{} systems across three build idioms; the whole pipeline ran on each and the
report was scored against the label. A true positive is one drift-report finding (a vanished or
new target, edge, rename, or violation) that the label expects, so a removed target yields several
findings. Precision \rqSixPrecision{}, recall \rqSixRecall{}, \rqSixFpPerPr{} false findings per
simulated pull request, rename accuracy \rqSixRenameAccuracy{}, and a one-line alias pin silences
the drift in \rqSixAliasSurvival{} pinned reruns. Splits and file moves between targets are not
among the operators.

\paragraph{Evidence degradation alone.} On the \rqSixDegradationSystems{} systems scored before
Squidex joined, inducing degradation (semantic toolchain lost; 10\,\% and 30\,\% of targets starved
of L0 evidence; build-graph extractor lost) produces \rqSixWouldBeFalse{} would-be-false
``disappeared'' findings, of which the stopping rules reclassify \rqSixSuppressionRate{} as
labelled coverage gaps; a coverage-naive differ reports every one as change (exact McNemar
$p\rqSixMcNemarP$).

\paragraph{Change and evidence loss at once (Table~\ref{tab:rq6b}).} The \rqSixCombMutations{}
mutations were replayed on \rqSixCombSystems{} systems with each degradation mode induced on the
mutated run, plus an adversarial mode that starves exactly the targets the mutation touches; each
of the \rqSixCombFindings{} expected findings was classified as named drift, labelled coverage gap,
visible only through the run's advisory downgrade, or silent. Under targeted starvation
\rqSixCombStarveMutatedDrift{} are drift, \rqSixCombStarveMutatedGap{} gaps, and
\rqSixCombStarveMutatedAdvisory{} advisory only. With the build-graph extractor lost entirely, the
first run reported an added or renamed target nowhere and was not even advisory
(\rqSixCombPrefixSilentTotal{} silent findings, the new targets of
\texttt{\rqSixCombPrefixSilentOps{}} mutations), because an empty first-party target set counted as
fully covered. The convention was changed (nothing covered, advisory) and the arm re-run:
\rqSixCombSilentTotal{} silent findings remain, the \rqSixCombLoseLZeroAdvisory{} former ones are
advisory, and no mode produced a false finding. The artifact keeps both runs
(\texttt{combined-prefix.json} is the pre-fix one).

\paragraph{Real history (Table~\ref{tab:rq6hist}).} \rqSixHistPairs{} consecutive release pairs of
\rqSixHistSystems{} systems, fixed in advance rather than selected by changelog, under the
committed rules held constant across each pair: \rqSixHistCommits{} commits,
\rqSixHistBuildCommits{} touching build files, \rqSixHistFindings{} findings,
\rqSixHistCoverageGaps{} coverage gaps; \rqSixHistPairsZero{} pairs with build-file churn produced
none. Of the \rqSixHistNewTargets{} first-party targets the releases added, the rules left
\rqSixHistNewUnmapped{} unmapped. Two raters labelled each finding independently (raw
agreement \rqSixHistLabelsAgree{}, pooled Cohen's $\kappa$ \rqSixHistLabelsKappa{}; both called
\rqSixHistArchitectural{} architectural, \rqSixHistNonArchitectural{} non-architectural,
\rqSixHistNoise{} extraction noise; the one disagreement is an added code-generator project
that one rater read as tooling rather than product architecture); the rater runbook and
worksheets are in the artifact.

\section{Performance per system}
\label{sec:tr-perf}
%% GENERATED by analysis/make_tables.py from results.csv -- DO NOT EDIT BY HAND.
\begin{table}[t]
  \centering
  \caption{RQ6: end-to-end wall-clock (s), cold vs.\ warm (incremental cache), speedup, and peak RSS (MB) of the whole process tree, ordered by first-party KLOC. Budgets: cold $\le$30\,min, warm $\le$3\,min. $^\ddagger$Anonymized proprietary system, non-replicable, corroborating only.}
  \label{tab:rq7}
  \footnotesize\setlength{\tabcolsep}{3.5pt}
  \begin{tabular}{@{}ll rr rrrr@{}}
    \toprule
    System & Lang & KLOC & Targets & Cold & Warm & Speedup & RSS \\
    \midrule
    toy & C\# & 0.1 & 4 & 15.0 & 4.2 & 3.6$\times$ & 258 \\
    libxml2 & C & 0.2 & 14 & 14.5 & 5.6 & 2.6$\times$ & 85 \\
    serilog & C\# & 7.1 & 1 & 62.2 & 13.7 & 4.5$\times$ & 516 \\
    eShop & C\# & 25 & 19 & 123.7 & 31.8 & 3.9$\times$ & 654 \\
    bash & C & 55 & 49 & 27.3 & 5.6 & 4.9$\times$ & 108 \\
    terminal & C++ & 83 & 43 & 330.1 & 34.3 & 9.6$\times$ & 93 \\
    P-1$^\ddagger$ & C++ & 124 & 52 & 676.7 & 205.5 & 3.3$\times$ & -- \\
    OpenCV & C++ & 132 & 24 & 696.6 & 69.1 & 10.1$\times$ & 697 \\
    abseil & C++ & 160 & 190 & 117.0 & 40.9 & 2.9$\times$ & 219 \\
    nopCommerce & C\# & 513 & 37 & 434.4 & 108.2 & 4.0$\times$ & 1018 \\
    ITK & C++ & 619 & 79 & 24600.2 & 546.8 & 45.0$\times$ & 1655 \\
    OrchardCore & C\# & 825 & 207 & 1022.8 & 838.0 & 1.2$\times$ & 1260 \\
    fmt & C++ & -- & 24 & 38.8 & 6.2 & 6.3$\times$ & 129 \\
    \bottomrule
  \end{tabular}
\end{table}


Table~\ref{tab:rq7} gives the cold and warm end-to-end wall-clock and the peak RSS of the whole
process tree per system, ordered by first-party KLOC, that the paper summarises: cold within the
30-minute budget on \rqSevenColdBudget{} of \rqSevenSystems{} systems, warm within 3 minutes on
\rqSevenWarmBudget{}, median speedup \rqSevenMedianSpeedup$\times$, log--log slope
\rqSevenScalingSlope{} against KLOC, peak RSS \rqSevenMaxPeakRssGb{}\,GB. ITK's cold run takes
\rqSevenItkColdMin{} minutes, \rqSevenItkDoxygenSharePct{}\,\% of it in Doxygen's XML pass,
against \rqSevenItkWarmMin{} warm; the L0-only run governance needs would take about
\rqSevenItkLZeroImpliedMin{} minutes (an estimate from the per-stage timings, not a measurement).

\section{Blind curation: rater agreement and effort}
\label{sec:tr-curation}
%% GENERATED by analysis/make_tables.py from results.csv -- DO NOT EDIT BY HAND.
\begin{table}[t]
  \centering
  \caption{RQ7 blind curation (terms: \S\ref{sec:design}; all-common projection; best SAR = best of ACDC/WCA/LIMBO, \texttt{dir} coinciding with the floor here). R$_1$ vs.\ R$_2$ = the reference authors' agreement, an inter-rater baseline; Edge-F1 = container-edge F1; $\kappa$ = raw Cohen's $\kappa$ on rater labels; confl. = rater clusters cutting a reference boundary; the consensus row is a sensitivity re-score, the second-curator row the curator-variance arm.}
  \label{tab:rq9}
  \footnotesize\setlength{\tabcolsep}{2.2pt}
  \begin{tabular}{@{}lll rrl@{}}
    \toprule
    System & Row & Kind & MoJo & a2a & Edge-F1\,/\,$\kappa$ \\
    \midrule
    eShop & no-curation floor & det. & 33.3 & 80.8 & -- \\
     & best SAR (LIMBO) & det. & 26.7 & 85.9 & -- \\
     & \textbf{blind curator} & expr. & \textbf{86.7} & \textbf{95.8} & 0.63 \\
     & \quad second blind curator (R-B) & expr. & 66.7 & 90.8 & 0.56 \\
     & \quad curator vs.\ curator agreement & agr. & 80.6 & 93.9 & ARI\,0.75 \\
     & \quad vs.\ consensus & expr. & 75.0 & 92.0 & 0.55 \\
     & R$_1$ vs.\ R$_2$ agreement & agr. & 66.5 & 89.7 & $\kappa$\,0.00; 0\,confl. \\
    \midrule
    Squidex & no-curation floor & det. & 70.0 & 92.0 & -- \\
     & best SAR (LIMBO) & det. & 60.0 & 93.0 & -- \\
     & \textbf{blind curator} & expr. & \textbf{70.0} & \textbf{95.8} & 0.67 \\
     & R$_1$ vs.\ R$_2$ agreement & agr. & 54.5 & 86.3 & $\kappa$\,0.39; 0\,confl. \\
    \bottomrule
  \end{tabular}
\end{table}

%% GENERATED by analysis/make_tables.py from results.csv -- DO NOT EDIT BY HAND.
\begin{table}[t]
  \centering
  \caption{RQ7 curation effort (Min = self-reported wall-clock; Ops = editor operations from the edit trace; Rules/Groups = the frozen rules; Prop. = groups kept from the \cli~\texttt{propose} draft) and inter-rater agreement structure ($\kappa_{\mathrm{al}}$ = $\kappa$ after optimal label matching; ARI = adjusted Rand index; NMI = normalized mutual information; Confl. as in Table~\ref{tab:rq9}; Cons. = clusters in the raters' consensus).}
  \label{tab:rq9effort}
  \footnotesize\setlength{\tabcolsep}{3.5pt}
  \setlength{\tabcolsep}{1.9pt}
  \begin{tabular}{@{}l rrrr r rrrr rr@{}}
    \toprule
    & \multicolumn{5}{c}{Curation effort} & \multicolumn{6}{c}{Rater agreement} \\
    \cmidrule(lr){2-6}\cmidrule(lr){7-12}
    System & Min & Ops & Rules & Groups & Prop. & $\kappa$ & $\kappa_{\mathrm{al}}$ & ARI & NMI & Confl. & Cons. \\
    \midrule
    eShop & 43 & 43 & 63 & 11 & 0/11 & 0.00 & 0.70 & 0.53 & 0.87 & 0 & 14 \\
    \quad R-B & 120 & 36 & 65 & 13 & 0/13 & -- & -- & -- & -- & -- & -- \\
    Squidex & 210 & 46 & 59 & 12 & 0/12 & 0.39 & 0.59 & 0.23 & 0.77 & 0 & 12 \\
    \bottomrule
  \end{tabular}
\end{table}


\paragraph{Rater agreement.} The plan's raw-$\kappa$ flag fired on both systems
(\rqNineEshopKappa{}, \rqNineSquidexKappa{}). Raw $\kappa$ scores exact label match, and the eShop
raters shared \rqNineEshopLabelsShared{} labels, which forces $\kappa=0$ arithmetically.
Label-free, eShop's raters agree (ARI \rqNineEshopAri{}, NMI \rqNineEshopNmi{}, aligned $\kappa$
\rqNineEshopKappaAligned{}) and Squidex's disagree (ARI \rqNineSquidexAri{}, NMI
\rqNineSquidexNmi{}, aligned $\kappa$ \rqNineSquidexKappaAligned{}); on both systems not one rater
cluster cuts across a reference boundary (\rqNineEshopRaterConflicts{} and
\rqNineSquidexRaterConflicts{} conflicts), so the partitions nest and every disagreement is about
altitude. On eShop both raters, who never saw the committed reference, reach at least
\rqNineEshopRaterRefMoJoMin{} MoJoFM against it (aligned
$\kappa\ge$\rqNineEshopRaterRefKappaAlignedMin{}, NMI $\ge$\rqNineEshopRaterRefNmiMin{}). The
committed reference is the GT-doc transcription the plan named as the scored reference before any
curator score was computed; the raters' consensus (\rqNineSquidexConsensusClusters{} clusters on
Squidex; a coarser altitude on eShop) is a sensitivity row, not a second reference; the eShop
curator scores \rqNineEshopCuratorVsConsensusMoJo{} against it.

\paragraph{Effort (Table~\ref{tab:rq9effort}).} Self-reported wall-clock, editor operations from
the edit trace, rule lines, and groups per curation: \rqNineEshopMinutes{} minutes,
\rqNineEshopOps{} operations, \rqNineEshopRuleLines{} rule lines for eShop's first curator;
\rqNineSquidexMinutes{} minutes, \rqNineSquidexOps{} operations, \rqNineSquidexRuleLines{} lines
for Squidex; the second eShop curator (\rqNineEshopSecondCuratorCode{})
\rqNineEshopSecondCuratorMinutes{} minutes over two sittings, scoring
\rqNineEshopSecondCuratorMoJo{} MoJoFM, a2a \rqNineEshopSecondCuratorAtoA{}, edge-F1
\rqNineEshopSecondCuratorEdgeFone{}, \rqNineEshopCuratorMoJoSpread{} points below the first; the
two curators agree at aligned $\kappa$ \rqNineEshopCuratorPairKappaAligned{} (ARI
\rqNineEshopCuratorPairAri{}, pairwise MoJoFM \rqNineEshopCuratorPairMoJo{}). No curator kept a
group from the \cli~\texttt{propose} draft (\rqNineEshopProposeDerivedPct\% and
\rqNineSquidexProposeDerivedPct\%). The edit traces, attestations, and frozen rules are tracked in
the artifact under \texttt{references/<system>/blind-curation/}.

\section{Fidelity tables of the paper}
\label{sec:papertables}
%% GENERATED by analysis/make_tables.py from results.csv -- DO NOT EDIT BY HAND.
\begin{table*}[t]
  \centering
  \caption{Corpus, RQ1 fidelity, and the A1/A5 ablation columns (terms: \S\ref{sec:design}). \#Cl./\#Ent./\#Tgt.\ = reference clusters, scored entities, first-party targets; Best = best all-common classical baseline by MoJoFM; A1 = container edges at L0 $\to$ all layers; A5 = container edges the never-drop-declared rule saves at threshold 3 / 8. GT-diag = inception-era diagram, architect-reconciled (P-1 curated with that diagram in view); $^\ddagger$anonymized proprietary system, non-replicable, corroborating only (entity count bucketed). P-1 floor/propose: its A2 rungs. $^\S$Squidex's only curated model is its reference-blind RQ7 curation.}
  \label{tab:corpus}\label{tab:rq1}\label{tab:rq3}
  \footnotesize\setlength{\tabcolsep}{3pt}
  \begin{tabular}{@{}lllll rrr rr rr rr r rr@{}}
    \toprule
    & & & & & & & & \multicolumn{2}{c}{Floor} & \multicolumn{2}{c}{Propose} & \multicolumn{2}{c}{Curated} & Best & A1 & A5 \\
    \cmidrule(lr){9-10}\cmidrule(lr){11-12}\cmidrule(lr){13-14}
    System & Lang & Build & Ref & Cur. & \#Cl. & \#Ent. & \#Tgt. & a2a & MoJo & a2a & MoJo & a2a & MoJo & MoJo & edges & t3/t8 \\
    \midrule
    ITK & C++ & CMake & GT-pub & assisted & 13 & 7657 & 79 & 84.2 & 96.9 & 97.8 & 98.3 & 99.6 & 98.4 & 93.9 & 17$\to$23 & 0/5 \\
    abseil & C++ & CMake & GT-exp & assisted & 21 & 809 & 190 & 85.1 & 89.9 & 100.0 & 100.0 & 98.8 & 95.2 & 97.8 & 79$\to$83 & 0/45 \\
    OpenCV & C++ & CMake & GT-doc & assisted & 14 & 1087 & 24 & 100.0 & 100.0 & 87.0 & 48.3 & 100.0 & 100.0 & 92.1 & 9$\to$9 & 0/5 \\
    bash & C & Makefile.in & GT-pub & assisted & 14 & 373 & 49 & 97.7 & 91.9 & 95.3 & 82.7 & 98.1 & 93.2 & 91.9 & 8$\to$17 & 0/5 \\
    libxml2 & C & automake & GT-pub & none & 18 & 82 & 14 & 76.0 & 11.4 & 73.5 & 5.7 & 76.0 & 11.4 & 31.4 & 3$\to$3 & 0/3 \\
    eShop & C\# & .NET SDK & GT-doc & sighted & 9 & 19 & 19 & 80.8 & 33.3 & 73.3 & 0.0 & 88.9 & 66.7 & 33.3 & 30$\to$30 & 0/30 \\
    OrchardCore & C\# & .NET SDK & GT-doc & sighted & 34 & 207 & 207 & 67.6 & 11.6 & 91.1 & 69.5 & 91.7 & 72.1 & 47.4 & 70$\to$72 & 0/26 \\
    nopCommerce & C\# & .NET SDK & GT-doc & sighted & 17 & 37 & 37 & 80.2 & 37.5 & 89.6 & 71.9 & 90.2 & 71.9 & 46.9 & 16$\to$30 & 0/0 \\
    Squidex$^\S$ & C\# & .NET SDK & GT-exp & blind & 12 & 15 & 26 & 92.0 & 70.0 & 86.8 & 50.0 & 95.8 & 70.0 & 70.0 & -- & -- \\
    P-1$^\ddagger$ & C++ & -- & GT-diag & sighted & 28 & 50-100 & 52 & 93.7 & 96.9 & 81.2 & 24.5 & 100.0 & 100.0 & 91.3 & -- & -- \\
    \bottomrule
  \end{tabular}
\end{table*}

%% GENERATED by analysis/make_tables.py from results.csv -- DO NOT EDIT BY HAND.
\begin{table*}[t]
  \centering
  \caption{RQ2: MoJoFM (top) and a2a (bottom) vs.\ the reference on the all-common projection, best per row in bold. Strict wins against the five all-common baselines: curated 6 of 9 systems on MoJoFM and 7 of 9 on a2a; floor 2 of 9 and 2 of 9; propose 4 of 9 and 4 of 9. $^\dagger$ARC on its own covered subset. $^\S$LLM, mean $\pm$ sd of five reruns; -- = refused or not run. $^\ddagger$Anonymized proprietary system, non-replicable, corroborating only. P-1 floor/propose: its A2 rungs.}
  \label{tab:rq2}
  \footnotesize\setlength{\tabcolsep}{4pt}
  \begin{tabular}{@{}lrrrrrrrrrr@{}}
    \toprule
    System & Floor & Prop. & Curated & ACDC & ARC$^\dagger$ & WCA & LIMBO & dir & cc & LLM$^\S$ \\
    \midrule
    ITK & 96.9 & 98.3 & \textbf{98.4} & 70.3 & 49.2 & -- & -- & 93.9 & 48.4 & -- \\
    abseil & 89.9 & \textbf{100.0} & 95.2 & 41.1 & 25.1 & 31.2 & 35.1 & 97.8 & 19.5 & 98.8$\pm$0.2 \\
    OpenCV & \textbf{100.0} & 48.3 & \textbf{100.0} & 81.7 & 28.7 & 25.9 & 46.9 & 92.1 & 25.4 & \textbf{100.0}$\pm$0.0 \\
    bash & 91.9 & 82.7 & \textbf{93.2} & 60.5 & 32.3 & 52.8 & 51.6 & 91.9 & 62.5 & 84.2$\pm$1.6 \\
    libxml2 & 11.4 & 5.7 & 11.4 & 20.0 & 36.9 & 31.4 & 28.6 & 5.7 & 25.7 & \textbf{62.8}$\pm$1.1 \\
    eShop & 33.3 & 0.0 & 66.7 & 6.7 & 33.3 & 20.0 & 26.7 & 33.3 & 6.7 & \textbf{86.7}$\pm$4.2 \\
    OrchardCore & 11.6 & 69.5 & \textbf{72.1} & 24.2 & 40.8 & 32.6 & 47.4 & 11.6 & 24.2 & 30.2$\pm$1.3 \\
    nopCommerce & 37.5 & \textbf{71.9} & \textbf{71.9} & 28.1 & 56.2 & 46.9 & 34.4 & 37.5 & 28.1 & 60.6$\pm$3.2 \\
    Squidex & 70.0 & 50.0 & 70.0 & 0.0 & 50.0 & 50.0 & 60.0 & 70.0 & 0.0 & \textbf{76.0}$\pm$13.6 \\
    P-1$^\ddagger$ & 96.9 & 24.5 & \textbf{100.0} & 64.7 & 16.1 & 12.9 & 22.5 & 91.3 & 12.3 & -- \\
    \midrule
    System (a2a) & Floor & Prop. & Curated & ACDC & ARC$^\dagger$ & WCA & LIMBO & dir & cc & LLM$^\S$ \\
    \midrule
    ITK & 84.2 & 97.8 & \textbf{99.6} & 73.2 & 81.3 & -- & -- & 80.4 & 85.5 & -- \\
    abseil & 85.1 & \textbf{100.0} & 98.8 & 83.0 & 81.9 & 82.8 & 81.5 & 90.0 & 79.9 & 97.6$\pm$0.1 \\
    OpenCV & \textbf{100.0} & 87.0 & \textbf{100.0} & 86.3 & 82.3 & 81.2 & 84.7 & 86.6 & 81.4 & \textbf{100.0}$\pm$0.0 \\
    bash & 97.7 & 95.3 & \textbf{98.1} & 80.7 & 84.0 & 89.0 & 85.1 & 97.0 & 90.5 & 92.6$\pm$0.2 \\
    libxml2 & 76.0 & 73.5 & 76.0 & 76.6 & 87.0 & 84.2 & 82.0 & 73.5 & 78.0 & \textbf{91.3}$\pm$0.3 \\
    eShop & 80.8 & 73.3 & 88.9 & 75.9 & 89.4 & 86.0 & 85.9 & 80.8 & 75.9 & \textbf{97.9}$\pm$0.7 \\
    OrchardCore & 67.6 & 91.1 & \textbf{91.7} & 80.2 & 85.1 & 82.0 & 83.6 & 67.6 & 79.0 & 81.8$\pm$0.5 \\
    nopCommerce & 80.2 & 89.6 & 90.2 & 76.5 & 90.1 & \textbf{90.7} & 87.8 & 80.2 & 76.5 & 84.2$\pm$12.3 \\
    Squidex & 92.0 & 86.8 & 95.8 & 69.8 & 88.4 & 88.4 & 93.0 & 92.0 & 69.8 & \textbf{96.1}$\pm$1.4 \\
    P-1$^\ddagger$ & 93.7 & 81.2 & \textbf{100.0} & 88.5 & 79.5 & 78.4 & 81.3 & 87.3 & 78.2 & -- \\
    \bottomrule
  \end{tabular}
\end{table*}

%% GENERATED by analysis/make_tables.py from results.csv -- DO NOT EDIT BY HAND.
\begin{table*}[t]
  \centering
  \caption{Cluster count $k$ (left; Ref = reference clusters; -- = not recorded) and post-hoc ARI ($\times$100, right; best per row in bold; $^\S$mean $\pm$ sd of the LLM reruns) of every partition in Tables~\ref{tab:rq1} and~\ref{tab:rq2}. $^\dagger$ARC on its own covered subset. $^\ddagger$Anonymized proprietary system, non-replicable, corroborating only. P-1 floor/propose: its A2 rungs.}
  \label{tab:rq2k}
  \scriptsize\setlength{\tabcolsep}{2.5pt}
  \begin{tabular}{@{}l r r r r r r r r r r@{}}
    \toprule
    System & Ref & Floor & Prop. & Cur. & ACDC & ARC$^\dagger$ & WCA & LIMBO & dir & cc \\
    \midrule
    ITK & 13 & 78 & 8 & 8 & 1034 & -- & -- & -- & 212 & 46 \\
    abseil & 21 & 90 & 15 & 14 & 27 & -- & 20 & 20 & 31 & 1 \\
    OpenCV & 14 & 14 & 2 & 14 & 108 & -- & 8 & 13 & 86 & 3 \\
    bash & 14 & 16 & 6 & 10 & 54 & 12 & 13 & 12 & 11 & 7 \\
    libxml2 & 18 & 3 & 1 & 3 & 4 & 13 & 10 & 9 & 1 & 2 \\
    eShop & 9 & 19 & 1 & 14 & 2 & -- & 8 & 7 & 19 & 2 \\
    OrchardCore & 34 & 201 & 20 & 16 & 13 & -- & 34 & 33 & 201 & 1 \\
    nopCommerce & 17 & 37 & 8 & 9 & 1 & -- & 17 & 16 & 37 & 1 \\
    Squidex & 12 & 13 & 6 & 10 & 1 & 7 & 7 & 9 & 13 & 1 \\
    P-1$^\ddagger$ & 28 & 49 & 6 & 27 & 52 & 13 & 2 & 18 & 86 & 1 \\
    \bottomrule
  \end{tabular}\hfill
  \begin{tabular}{@{}l r r r r r r r r r r@{}}
    \toprule
    System & Floor & Prop. & Cur. & ACDC & ARC$^\dagger$ & WCA & LIMBO & dir & cc & LLM$^\S$ \\
    \midrule
    ITK & 22.6 & 84.6 & \textbf{99.3} & 2.9 & -- & -- & -- & 11.1 & -6.6 & -- \\
    abseil & 44.1 & \textbf{100.0} & 93.5 & 9.5 & -- & 11.4 & 4.2 & 62.8 & 0.0 & 85.8$\pm$0.1 \\
    OpenCV & \textbf{100.0} & 26.9 & \textbf{100.0} & 32.6 & -- & -0.7 & 4.9 & 40.8 & -0.2 & \textbf{100.0}$\pm$0.0 \\
    bash & 85.7 & 73.0 & \textbf{86.4} & 16.3 & 3.6 & 18.9 & 15.2 & 78.8 & 38.0 & 57.6$\pm$1.9 \\
    libxml2 & 1.6 & 0.0 & 1.6 & 0.8 & -2.1 & 16.8 & 10.0 & 0.0 & 8.0 & \textbf{46.1}$\pm$1.0 \\
    eShop & 0.0 & 0.0 & 45.1 & -1.4 & -- & -0.7 & 1.5 & 0.0 & -1.4 & \textbf{75.2}$\pm$6.8 \\
    OrchardCore & 0.0 & 59.4 & \textbf{84.1} & -1.8 & -- & 3.5 & 10.2 & 0.0 & 0.0 & 2.3$\pm$0.4 \\
    nopCommerce & 0.0 & \textbf{86.8} & 64.6 & 0.0 & -- & 16.4 & 8.1 & 0.0 & 0.0 & 23.7$\pm$1.7 \\
    Squidex & 0.0 & 20.7 & 20.9 & 0.0 & 10.6 & 7.5 & 17.5 & 0.0 & 0.0 & \textbf{44.0}$\pm$20.0 \\
    \bottomrule
  \end{tabular}
\end{table*}


Tables~\ref{tab:rq1}, \ref{tab:rq2}, and \ref{tab:rq2k} are the paper's Tables I--IV, reproduced
from the same generated files so that the captions above can refer to them. The per-system
fidelity CSV (\texttt{fidelity-by-system.csv}) in the artifact carries the same numbers with the
all-common entity counts.

\bibliographystyle{IEEEtran}
\bibliography{references}

\end{document}
